\section{Resource Model and Execution Contracts}
\label{sec:scope}

\subsection{Representations and one notation}
The target is a complete request or forward output under a declared model,
arithmetic, placement, and quality profile $\pi$. Text decoding is the current
instance: prefill, multiple decode steps, and KV/state lifetime. Vision,
audio, and combined modalities require their own complete-output contracts
and remain future implementation and evaluation work.

A representation $r$ identifies payload, decoding metadata, expansion, and
required physical copies. Checkpoint bytes $S_{\mathrm{ckpt}}$ include the
specified files and container layout. Complete compact weights
$W_{\mathrm{compact}}$ include all encoded weights, scales, and decoding
metadata. Full-resident weight demand $W_{\mathrm{full}}(r)$ is the footprint
of all weights under representation and placement $r$; it is weight-only,
not the complete inference working set. Expanded execution weights may be
larger than compact weights and need not all be live at once. Aliased views
are not extra copies; genuinely copied payloads are.

\begin{table}[tbp]
\centering
\small
\begin{tabularx}{\linewidth}{@{}l>{\raggedright\arraybackslash}X@{}}
\toprule
Symbol & Meaning and units \\
\midrule
$\pi,r$ & Execution/quality profile and weight representation/placement \\
$S_{\mathrm{ckpt}},W_{\mathrm{compact}},W_{\mathrm{full}}(r)$ & Checkpoint, complete compact, and full-resident weight bytes \\
$d\in\mathcal D_{\mathrm{phys}},B_d$ & Disjoint physical domain and available physical byte capacity \\
$a\in\mathcal D_{\mathrm{acct}},b_a$ & Covered accounting domain and configured limit in that domain's unit \\
$M_d(t),M_d^{\mathrm{peak}}$ & Deduplicated physical working-set bytes at time $t$, and its true peak \\
$\widehat M_d(t_i),\widehat M_d^{\mathrm{max}}$ & Reported sample and maximum over stated sampling times; not a true peak \\
$U_a(t),c_a$ & Charged occupancy and additional proposed claim, in domain $a$'s unit \\
$\delta,\rho_a,b_a^{\mathrm{eff}}(\delta)$ & Demand class, execution reserve, and its policy-accessible accounting ceiling \\
$F_a,R_a,I_a,L_a,T_a$ & Free, reserved, provider-in-flight, live, and retained claims \\
$P=(K,\boldsymbol q,\boldsymbol\ell)$ & Frozen plan: exact key, stage claim vectors, resident claim vector \\
$K,g_s,g_x$ & Full resource key, source generation, destination reuse generation \\
$o,e,J$ & Owned operation, completion evidence, provisional owner journal \\
$\Omega,k,w,\mu_o[s]$ & Post-detachment worklist, per-pass visit limit, persistent cursor, and exact-stage cancellation marker \\
$\mathcal V_q,L_q,H_q,\Delta_q$ & Request $q$'s invocations, logical reader bytes, facade H2D bytes, and $H_q-L_q$ \\
$N_j,I_j,E_j$ & Expert entries at invocation boundary $j$, installs/uploads, and evictions \\
$X_j,Q_j,S_j,C_j^{\mathrm{cache}},Z_j$ & Resident, pending-upload, scratch, cache charge, and retained Unknown subset gauges \\
$C_{\mathrm{peak}}^{\mathrm{cache}}$ & Lifetime cache-accounting high-water; not a total physical-owner peak \\
$P_q^{\mathrm{net}},P_q^{\mathrm{blk}}$ & Physical network and backing-block bytes attributed to request $q$ \\
$b_{\mathrm{cache}},b_{\mathrm{scratch}}$ & Configured cache-credit and scratch limits; neither is a physical-domain bound \\
\bottomrule
\end{tabularx}
\caption{Notation used by the contracts, algorithms, and evaluation.
Accounting limits are never substituted for physical capacity.}
\label{tab:symbols}
\end{table}

\subsection{Physical capacity and measured working set}
Physical capacity counts disjoint resident-memory pools once:
\begin{equation}
 B_{\mathrm{mem}}=\sum_{d\in\mathcal D_{\mathrm{phys}}}B_d.
 \label{eq:physical-capacity}
\end{equation}
For a declared representation and placement, consider
\begin{equation}
 W_{\mathrm{full}}(r)>B_{\mathrm{mem}}.
 \label{eq:full-residency}
\end{equation}
The inequality rules out aggregate full residency of \emph{that}
representation/placement, not streaming feasibility. Disk is excluded.
Discrete host and GPU pools may be added; host and accelerator aliases of
one physical pool may not. Unified virtual addressing over separate physical
pools does not merge those pools. Report installed capacity, fixed exclusions,
co-tenants, copies, and placement/migration policy for each evaluated domain.
Per-domain accessibility and placement remain constraints even when the
aggregate sum is sufficient.

Define $M_d(t)$ as physical bytes resident in domain $d$ that are attributable
to the selected inference deployment at time $t$. Count resident weights,
expansion/workspace, activations, model state, staging, outstanding and retained
buffers, metadata, allocator/driver reserves, and attributable file-backed
pages, deduplicating shared physical allocations. A full attribution requires
fixed and external consumers to be measured or excluded explicitly; a
process heap or device counter alone does not supply it. The desired
feasibility condition and the reported sampling statistic are different:
\begin{equation}
\begin{aligned}
 M_d^{\mathrm{peak}}&=\sup_{t\in\mathcal T}M_d(t)\leq B_d
       &&\text{for each }d,\\
 \widehat M_d^{\mathrm{max}}&=\max_{t_i\in\mathcal T}
       \widehat M_d(t_i)
       &&\text{over the declared samples.}
\end{aligned}
\label{eq:working-set}
\end{equation}
$\mathcal T$ is the declared startup/request/cleanup interval. A sampled
maximum can miss a transient peak; selected-device telemetry also includes
unattributed processes, and tree PSS does not cover every machine consumer.
Thus those proxies are not silently identified with $M_d(t)$. Independent
component maxima must not be added into a simultaneous working-set peak.
RQ1 uses Table~\ref{tab:tp1-memory}; the storage-demand figures in
Section~\ref{sec:rq2} diagnose cache and transfer counters without turning
sample maxima into certified physical bounds.

The storage-necessity gate additionally asks whether \emph{all compact
weights and metadata} can remain resident with online expansion under the
quality contract. Failure to fit expanded weights is insufficient if that
compact alternative fits. A low configured cache budget is likewise not
physical host-memory insufficiency. Storage necessity is scoped to evaluated
representations, not a proof against every possible compression scheme.

\subsection{Accounting and admission contract}
A covered accounting domain $a$ may count bytes, slots, pages, or queue entries.
Its limit $b_a$ is configured, not installed physical capacity; accounting
domains can overlap or omit parts of a pool and cannot in general be summed.
A cache limit covers its cache, not the whole working set. The necessary
covered-domain capacity condition is
\begin{equation}
 \forall a\in\mathcal D_{\mathrm{acct}},\qquad
 U_a(t)+c_a\leq b_a.
 \label{eq:admission}
\end{equation}
Here $c_a$ is the additional charge to acquire, not a second charge for an
already held grant. A revisit preserves that grant and rechecks only unmet
claims. The inequality is not sufficient for admission: the owner must also
validate canonical claims and units, arithmetic overflow,
source/placement/backend compatibility,
actual reservation and grant authority, and the selected demand policy.
With execution reserve $0\leq\rho_a\leq b_a$, the policy-accessible ceiling is
\begin{equation}
 b_a^{\mathrm{eff}}(\delta)=
 \begin{cases}
 b_a,&\delta=\mathsf{Required},\\
 b_a-\rho_a,&\delta\in\{\mathsf{Warmup},\mathsf{Prefetch}\}.
 \end{cases}
 \label{eq:execution-reserve}
\end{equation}
A non-required demand cannot consume the execution reserve even if
Equation~\ref{eq:admission} holds. Required execution may use that reserve,
not exceed the hard ceiling. The reserve protects accounting capacity from
speculation, not eventual progress against a stalled provider.

The abstract admission result has three cases. \textsf{Rejected} includes
invalid plans and an intrinsically impossible stage: for example, a required
claim exceeding $b_a$, or a non-required claim exceeding its unchanged
policy-accessible ceiling even with no other occupancy. Waiting cannot fix
such a stage; it needs rejection or an explicitly changed, revalidated plan
or demand policy. \textsf{Wait} means the otherwise admissible stage is
blocked by current occupancy, including reserve-protected capacity. It retains
existing claims and registers a revisit condition rather than recharging or
redispatching them. \textsf{Admitted} requires all authority, compatibility,
and policy checks as well as capacity. These are explanatory result classes,
not a new backend API, and do not imply that all future stages can be admitted
atomically or that their resource waits are acyclic.

Each claim has an authoritative owner, unit, and end condition. Canonicalizing
same-kind claims is not deduplication of arbitrary aliased allocations. A
bounded request channel does not bound all downstream buffers. Map claims to
the actual ledger and physical allocations on the selected path; do not infer
one global broker from a common interface.

The plan $P$ binds $K$, read/upload/install claim vectors $\boldsymbol q$,
and persistent residency $\boldsymbol\ell$. $K$ includes model instance,
source identity, semantic content descriptor, encoding, backend/device,
and $g_s,g_x$. Alignment and multiple extents may change staged byte counts;
expansion may change execution bytes. The destination authority supplies the
actual reservation and $g_x$ before hard admission. Publication consumes
validated bindings, not requested addresses. Algorithm~\ref{alg:normal-materialization}
uses the capacity and policy checks in Equations~\ref{eq:admission}
and~\ref{eq:execution-reserve} before each physical dispatch.

OS page cache, storage-server cache, metadata, and driver reserves are not
implicitly under runtime credit control. Resident file-backed pages count as
memory; the backing file adds no capacity. Swap must be disabled or explicitly
accounted for. A logical read, page fault, or migration between physical pools
is not evidence of a physical storage read. RQ2 consequently distinguishes
$L_q,H_q,\Delta_q,P_q^{\mathrm{net}},P_q^{\mathrm{blk}}$ rather than translating
them into one bandwidth number.

\subsection{Quality, completion, and progress contracts}
Protocol exactness means compatible sharing, exact completion identity,
authorized publication, and no reuse before physical references end. It is
separate from numerical/task quality. A quality contract fixes complete
outputs, independent reference, arithmetic/state precision, inputs, and
acceptance thresholds. Visible text, usage, generated IDs, full-vocabulary
logits, ordered routes, and internal states are different checks; argmax
agreement does not establish numerical agreement. Quantized checkpoint
storage alone does not specify kernel or KV arithmetic.

Completion evidence must name its operation, full key, submitted stage,
generations, outcome-specific byte count, and any required physical fence.
Rejecting an event with mismatched identity, stage, generation, or byte
contract is not the same as accepting a matching event whose validated
\texttt{CompletionOutcome::Stale} is a non-success result. The former supplies
no mutation/release authority; the latter selects terminal cleanup of its
named operation, never readiness or publication of a newer generation.
Returning named read custody is separate from releasing the read-slot and
storage-read claims; pinned host capacity remains charged through upload or
proven cleanup. Publication preflight similarly distinguishes Ready, Blocked
by live replacement custody, and Rejected by an invalid binding/conflict.
Rejected publication retains its journal and remaining custody for cleanup.

Only validated producer success proves readiness of that stage, not
quiescence of a later consumer. Only owner-accepted cancellation selects a
canceled caller outcome and detaches demand; a cancellation request or provider acceptance
alone does not prove physical completion. Shared loads retain their remaining
owners. Pending/Idle is retained as Pending only for a positively known
original pollable obligation; otherwise it retains Unknown custody. Unknown
means remaining physical references cannot be proved. Unknown retains allocation
and credits, possibly indefinitely. A late exact scoped completion can close
its named registry obligation but cannot clear physical-owner poison or
quarantine. Bounded cleanup passes retain unresolved obligations and a
traversal cursor across wakes and revisits, so the same first items do not hide unvisited
obligations. Ambiguous cancellation replies remain non-retryable without
exact owner evidence; they are not treated as known non-acceptance. Neither
that scheduling contract nor Unknown retention promises eventual reclamation
or availability.

Conditional liveness requires feasible per-domain resources, no cyclic
resource waits, fair progress and completion handling, live providers, and
eventually observable completion. A blocking provider or stalled progress
worker can violate these conditions. An asynchronous API, event-driven
progress, and hardware overlap are separate properties. Neither algorithms
nor current receipts establish global nonblocking execution, lock freedom,
durable crash recovery, or arbitrary multi-owner atomic commit.
